\section{A uniform isogeny bound for two families of sextic fields}
\label{sec:sextic}

Let $E$ be an elliptic curve over a number field $K$. A cyclic isogeny of
prime degree $p$ over $K$ gives a Galois-stable line in $E[p]$. Write
$\lambda$ for the character describing the action on the line, and $\chi_p$
for the mod-$p$ cyclotomic character. We use \emph{Momose type 2} in the
sense of Banwait and Derickx, with $\lambda^{12}=\chi_p^6$ and
$p\equiv3\pmod4$ as in their type 2 sieve. The bounds below concern this
particular type. A bound independent of the discriminant of $K$ is the
relevant uniformity question.

\paragraph{Sextic-field theorem.}
Suppose the normal closure $L$ of $K/\mathbb Q$ acts on the six embeddings
of $K$ as either
$A_4$ on the six edges of a tetrahedron or $A_5$ on its six Sylow
$5$-subgroups. Then every odd Momose type 2 isogeny prime belongs to an
explicit set of $66$ primes, with largest member $253{,}507$. The set
is exactly the set compatible with the local obstruction below for a single field in each family simultaneously, among primes $p\equiv3\pmod4$. Whether the largest prime occurs as
the degree of an actual type 2 isogeny over one of these fields is a separate
existence question. The $A_5$ family consists of primitive sextic fields,
since the stabilizer of a Sylow $5$-subgroup is maximal in $A_5$ (the action
is doubly transitive, as the stabilizer's $5$-cycle permutes the other
five points).

\paragraph{Two inputs.}
Banwait and Derickx, \href{https://doi.org/10.2140/ant.2025.19.1147}
{\emph{Explicit isogenies of prime degree over number fields}},
Proposition 6.1, gives the following necessary condition. For any
rational prime $q$ splitting in $\mathbb Q(\sqrt{-p})$ and any prime
above it in $K$ of odd residue degree $f$,
\begin{equation}
 4q^f\ge p \quad\text{or}\quad p\mid q^{2f}+q^f+1.
 \tag{CC}\label{eq:cc}
\end{equation}
The second alternative matters for potentially multiplicative reduction.
There is no assumption that $q$ be unramified in $K$.
The same authors, in \href{https://arxiv.org/pdf/2302.08350}
{\emph{Towards strong uniformity for isogenies of prime degree}},
Appendix A, Algorithm A.8 and Example A.13, report the endpoint
$253{,}507$ for the arithmetic test using residue degrees $1,3$. They
apply this test to cubic fields and ask how their method might treat
fields of even degree. Section~\ref{sec:strict-cc} verifies the needed
strict version: every prime $p\equiv3\pmod4$ larger than $253{,}507$
has a prime $q$ splitting in $\mathbb Q(\sqrt{-p})$ with
\[
 4q^3<p,\qquad p\nmid(q^2+q+1)(q^6+q^3+1).
\]
The small split prime bound uses Bach and Sorenson for
$\mathbb Q(\sqrt{-p})$. Its zeta function is
$\zeta(s)L(s,\chi_{-p})$. This explains why Dirichlet GRH
supplies the analytic hypothesis even when $L/\mathbb Q$ is nonabelian.
In using the result one must retain the strict inequality even when
$p$ lies just below $4q^3$. The endpoint $p=3$ has no $q$ with $4q<p$
and so passes the arithmetic test vacuously.

\paragraph{Residue degrees, including ramification.}
Write $G=\operatorname{Gal}(L/\mathbb Q)$, $H=\operatorname{Gal}(L/K)$
and $X=G/H$. At a prime of $L$ over $q$, let $I\triangleleft D$ be the
inertia and decomposition groups. Take a lift $g\in D$ of the Frobenius
generator of $D/I$. The $D$-orbits on $X$ correspond to primes of $K$
over $q$. In each such orbit the inertia orbits are permuted transitively
by $g$, and their number is the residue degree $f$.
Indeed, for a point with stabilizer $D_x$ the number is $[D:I D_x]$,
the residue degree in the associated local subfield. Consequently $f$
divides the length of the $g$-cycle of any point in that orbit.
An odd cycle of length $1$ or $3$ supplies a prime with $f=1$ or $3$.

\begin{center}
\begin{tabular}{lrrrr}
 & $1^6$ & $1^2 2^2$ & $3^2$ & $1\,5$\\
\hline
$A_4$ on edges & 1 & 3 & 8 & 0\\
$A_5$ on Sylow $5$-subgroups & 1 & 15 & 20 & 24
\end{tabular}
\end{center}
Here the entries count group elements of the indicated cycle shape.
For $A_4$ the double transpositions fix two edges. For $A_5$ the
stabilizer of a point is a dihedral group of order $10$. Its five
involutions account, by counting incidences across the six stabilizers,
for two fixed points per involution. Each element of order five fixes
its unique Sylow subgroup. Each element of order three acts without a
fixed point. These observations prove the table and show why the same
arithmetic sieve applies to both sextic actions. Testing the primes
$p\equiv3\pmod4$ up to $253{,}507$ gives $66$ survivors. The complete list
is in \path{checks/isogeny_sieve.out}. This proves the upper bound.

\paragraph{General uniform bound.}
Suppose more generally that every element of a specified transitive
permutation group on $G/H$ has an odd cycle of length at most $m$. Then
each prime of $\mathbb Q$ has an odd residue degree $f\le m$ somewhere
above it in $K$. Write $U(p)=(4\log p+10)^2$. The explicit splitting
bound used in Banwait--Derickx, Appendix A, gives a split prime
$q\le U(p)$ in $\mathbb Q(\sqrt{-p})$. Thus every Momose type 2 prime
must satisfy
\[
 p\le \max\{4U(p)^m,\ U(p)^{2m}+U(p)^m+1\}.
\]
If both expressions were smaller than $p$, both alternatives in
\eqref{eq:cc} would fail for the prime of odd residue degree just found.
The inequality gives an effective bound depending only on $m$, since its
right side is a fixed polynomial in $\log p$.

A variant allows a fixed abelian extension $F/\mathbb Q$ inside $L$.
It suffices that every element of $\operatorname{Gal}(L/F)$ have a short
odd cycle. Apply the least split prime estimate to the abelian field
$J=F\,\mathbb Q(\sqrt{-p})$ instead. Its degree is at most
$2[F:\mathbb Q]$, and its absolute discriminant is at most
$|D_F|^2p^{[F:\mathbb Q]}$ by the compositum discriminant inequality.
The resulting $q\ll_F (1+\log p)^2$ splits in both fields. Its
decomposition group in $L$ acts trivially on $F$, so the same argument
works even if $q$ ramifies in $L/F$. For example, the edge action of
$S_4$ has this property with $m=3$ once its quadratic sign field is fixed.
This yields a bound uniform over these sextics with fixed sign field.

\paragraph{Sharpness of the local obstruction.}
Fix any one of the surviving primes $p$. Let $S$ consist of the primes
$q<p/4$ splitting in $\mathbb Q(\sqrt{-p})$. We first construct a cyclic
cubic field $C$ in which all primes in $S$ are inert. Consider
\[
 u_t(X)=X^3-tX^2-(t+3)X-1,\qquad
 \operatorname{disc}(u_t)=(t^2+3t+9)^2.
\]
For every prime $q$, the values $0,-1$ are never roots. Each of the other
$q-2$ residue classes $x$ gives at most one $t$ with $u_t(x)=0$,
since the coefficient of $t$ is $-x(x+1)$. Thus there is some $t$ with
$u_t$ irreducible modulo $q$. Choose these parameters simultaneously
by the Chinese remainder theorem, adding an auxiliary prime if $S$ is
empty. Then $u_t$ is irreducible over $\mathbb Q$, its splitting field
$C=\mathbb Q(\alpha)$ is cyclic cubic, and every $q\in S$ is unramified
and inert because its reduction is separable and irreducible.

Here are explicit details for extending $C$ to a tetrahedral normal
closure. Choose a rational prime splitting completely in $C$ and use
weak approximation to find $\beta\in C^\times$ with valuations $(1,0,0)$
at its three primes. For a generator $\sigma$ of
$\operatorname{Gal}(C/\mathbb Q)$ set $a=\beta/\sigma\beta$.
The classes of $a$ and $\sigma a$ are independent in
$C^\times/(C^\times)^2$, since their valuation vectors modulo two are
two distinct nonzero even-weight vectors on the three primes.
The relation $a\,\sigma a\,\sigma^2 a=1$ gives a Galois field
$L=C(\sqrt a,\sqrt{\sigma a})$ over $\mathbb Q$. Its order-$3$ quotient
permutes the three nontrivial square classes cyclically. An order-$3$
subgroup exists by Cauchy's theorem, so the group is the nontrivial
semidirect product $V_4\rtimes C_3\cong A_4$. The sextic subfield
$K=C(\sqrt a)$ has the tetrahedral edge action.

For $q\in S$, every residue degree in $K$ is a multiple of $3$ and at
most $6$. Thus the only odd degree to consider is $3$. If $4q^3<p$,
then $q^2+q+1<p$. Since $p$ passed the uniform sieve, we must have
$p\mid q^6+q^3+1$. If $4q^3\ge p$, the first clause in
\eqref{eq:cc} already holds. Every other split prime has $4q\ge p$.
This proves \eqref{eq:cc} for the constructed field at every prime,
establishing the asserted exact classification of compatible primes.

For the endpoint $p=253{,}507$, the set $S$ has $3{,}152$ primes, from
$41$ through $63{,}361$. Already $4\cdot41^3>p$.
For this witness the supplied parameter is rational, $t=A/B$ in lowest
terms. To shorten the defining polynomial, the checker first combines
admissible residues into $t_0$ modulo $M=17^2\prod_{q\in S}q$, with
$t_0\equiv0\pmod{17^2}$. The Euclidean algorithm gives a short pair
$(A,B)$ with $A\equiv t_0 B\pmod M$ and $\gcd(B,M)=1$.
The largest absolute coefficient of the integral polynomial
$B X^3-A X^2-(A+3B)X-B$ has $6{,}796$ decimal digits. A monic integral
polynomial for the same field is
$Y^3-A Y^2-B(A+3B)Y-B^3$, with largest coefficient of $20{,}382$
digits in absolute value. Since the
denominator is a unit at each prescribed prime, reduction and
irreducibility work as before. Modulo $17$, $u_t$ has three simple
roots $3,4,10$ and $u_t(3)\equiv17\pmod{17^2}$. Hence
$\beta=\alpha-3$ has precisely the valuations needed above. One can take
$\sigma(\alpha)=-1/(\alpha+1)$, as substitution into $u_t$ verifies.
The two integers in \path{checks/sextic_parameter.txt} define a specific sextic-field
witness via two short equations. It also gives a cyclic cubic witness
to the local compatibility question discussed in Remark A.14 of the same
preprint. Their cubic example has coefficients exceeding $27{,}380$
decimal digits and enforces inertness at every prime below $p/4$.
The weaker splitting prescriptions used here suffice for the same
necessary condition.
